
\section*{Глава 3. Реализация веб-приложения}
\addcontentsline{toc}{section}{Глава 3. Реализация веб-приложения}
\stepcounter{section}

\subsection{Инициализация и развёртывание}

Приложение построено по паттерну \textbf{Application Factory}: функция
\texttt{create\_app()} создаёт экземпляр Flask, подключает \texttt{SQLAlchemy}
и \texttt{Flask-Migrate}, регистрирует пять блупринтов (\texttt{main},
\texttt{auth}, \texttt{games}, \texttt{users}, \texttt{notifications}) и
устанавливает контекстный процессор, подсчитывающий количество непрочитанных
уведомлений для текущего пользователя при каждом запросе.

Конфигурация передаётся через переменные окружения (\texttt{.env}): строка
подключения к PostgreSQL, секретный ключ сессии и корневой путь файлового
хранилища. Такой подход позволяет запускать приложение в разных средах без
изменения исходного кода.

Среда выполнения описана в \texttt{docker-compose.yml}~\cite{docker_docs} и включает
три сервиса: \textbf{db} (PostgreSQL~15~\cite{postgresql_docs}, данные сохраняются в
именованном томе), \textbf{web} (Flask + Gunicorn, Python-код запекается в образ при
сборке) и \textbf{nginx}~\cite{nginx_docs} (обратный прокси, отдаёт статические файлы
и игровые сборки напрямую из файловой системы, не нагружая Flask-процесс).

Миграции схемы управляются через \texttt{Flask-Migrate} (Alembic); директория
\texttt{migrations/} примонтирована как том, поэтому \texttt{flask db upgrade}
выполняется внутри работающего контейнера без пересборки образа.

\subsection{Аутентификация и ролевая модель}

Регистрация и вход реализованы в блупринте \texttt{auth}. Пароль хешируется
функцией \texttt{werkzeug.security.generate\_password\_hash} (pbkdf2-sha256)
и не хранится в открытом виде. После успешного входа идентификатор,
имя, роль и URL аватара записываются в подписанную серверную сессию Flask.

Платформа использует четыре роли. \textbf{Player} --- базовая роль при
регистрации: каталог, запуск игр, комментарии. \textbf{Developer} ---
выдаётся мгновенно по запросу пользователя в профиле, открывает форму
публикации игр. \textbf{Moderator} --- при запросе все действующие
модераторы и администраторы получают уведомление; роль назначается только
после явного одобрения администратором. \textbf{Admin} --- назначается вручную,
имеет полный доступ.

Контроль доступа в маршрутах реализован декоратором \texttt{@login\_required}
и проверками \texttt{session.get('role')}; при нарушении прав Flask возвращает
код~403.

\subsection{Публикация игр и каталог}

Публикация доступна пользователям с ролью \texttt{developer} или \texttt{admin}
через форму \texttt{/games/publish}. Разработчик задаёт название, описание,
жанры, цвет фона страницы (\texttt{bg\_color}), ссылку на обложку и ссылки
на социальные сети. URL-идентификатор (\texttt{slug}) генерируется
транслитерацией из названия и гарантирует уникальность.

Для запуска игры в браузере поддерживаются два способа. Первый ---
\textbf{внешняя ссылка}: разработчик вводит готовый URL (itch.io или собственный
хостинг), платформа встраивает его в \texttt{<iframe>}. Второй ---
\textbf{ZIP-загрузка}: архив Unity WebGL распаковывается в директорию
\texttt{storage/games/<slug>/}, найденный \texttt{index.html} записывается
в поле \texttt{iframe\_url} и далее обслуживается статически через nginx.

При публикации всем активным модераторам и администраторам отправляется
уведомление с предложением проверить контент. Модератор может снять игру с
публикации прямо со страницы игры; автор при этом получает уведомление типа
\texttt{game\_removed}.

Каталог (\texttt{/games/}) поддерживает полнотекстовый поиск по названию и
фильтрацию по жанру; сортировка --- по дате и количеству запусков.

\subsection{Страница игры и запуск WebGL}

Страница игры состоит из hero-секции и полноширинного \texttt{<iframe>}.
Цвет фона hero-секции берётся из \texttt{bg\_color}; свойство модели
\texttt{bg\_is\_dark} вычисляет относительную яркость цвета и автоматически
выбирает читаемый цвет текста --- светлый или тёмный.

Unity WebGL-сборки~\cite{unity_webgl} поставляются в сжатом виде (Brotli или Gzip).
Без явного указания \texttt{Content-Encoding} браузер не может распаковать файлы
и игра не запускается. В \texttt{nginx.conf} для каждого расширения
(\texttt{.wasm.br}, \texttt{.js.br}, \texttt{.data.br} и их \texttt{.gz}-аналогов)
прописан отдельный \texttt{location}-блок с заголовками:
\texttt{Content-Encoding}, корректным \texttt{Content-Type}, а также
\texttt{Cross-Origin-Opener-Policy: same-origin} и
\texttt{Cross-Origin-Embedder-Policy: require-corp}~\cite{mdn_coep} --- они
обязательны для активации \texttt{SharedArrayBuffer} в современных браузерах.

Заголовки COOP/COEP выставляются только на location \texttt{/storage/games/},
но не на саму страницу \texttt{/games/<slug>}, чтобы не блокировать загрузку
внешних обложек с других доменов.

\subsection{Комментарии, оценки и уведомления}

Оценка игры совмещена с комментарием: при отправке формы пользователь
опционально выбирает от 1 до 5 звёзд. Значение записывается в поле
\texttt{rating} таблицы \texttt{comments}. Средний рейтинг вычисляется
SQL-агрегацией (\texttt{func.avg()}) непосредственно в базе данных и
округляется до одного десятичного знака; результат отображается в hero-секции.

Система уведомлений полностью хранится в базе данных: при наступлении события
сервер вставляет строки в таблицу \texttt{notifications} для каждого получателя.
Контекстный процессор подсчитывает количество непрочитанных записей и передаёт
значение во все шаблоны --- оно отображается счётчиком на иконке колокольчика.
При открытии страницы \texttt{/notifications} все записи помечаются прочитанными
одним \texttt{UPDATE}-запросом.

Профиль пользователя позволяет сменить роль (player~→~developer мгновенно;
player / developer~→~moderator через заявку) и загрузить аватар. Загруженный
файл сохраняется в \texttt{storage/avatars/} под уникальным именем, URL
обновляется в базе данных и в текущей сессии.

\subsection{Доступность исходного кода}

Исходный код платформы WhyGame опубликован в открытом репозитории на GitHub
и доступен по адресу:

\begin{center}
    \url{https://github.com/whynotfu/WhyGame}
\end{center}

Репозиторий содержит Flask-приложение, конфигурацию Docker Compose, настройки
nginx, Alembic-миграции и LaTeX-документацию. Файлы \texttt{storage/}
(игровые сборки и аватары) исключены из репозитория через \texttt{.gitignore},
поскольку являются пользовательским контентом и не должны версионироваться.
